% !PRINTED 282
\weberpage{282}
%<<CONT
\[
\begin{array}{ll}
\text{for } n=7, & m=24,\\
\textquotedblbase\quad n=23, & m=24,\\
\textquotedblbase\quad n=31, & m=96,\\
\textquotedblbase\quad n=47, & m=48,\\
\textquotedblbase\quad n=71, & m=72,\\
\textquotedblbase\quad n=15, & m=72.
\end{array}
\]

In the cases 23, 47, 71 the degree of the correspondence is therefore
equal to the degree of the corresponding modular equation, that is,
equal to the number of the transformations in question; in the cases
$n=7, 31, 15$ the degree is \gesperrt{three times} the degree of the
modular equations.

What do these supernumerary points signify? They are explained by
the fact that in these cases, in which $n\equiv 0, 1$
$\pmod 3$, the exponents $h+2k$ in the function $\Phi_{a,c,\partial}$ [§ 75, (5)], which by the
above equation (5) is equal to a constant, leave the same residue
with respect to the modulus 3 in all terms, and
that accordingly the equation $\Phi=0$, and likewise the equations (2),
remain satisfied if $\xi, \eta, \zeta$ are multiplied by an arbitrary third root of unity. Any three points of the
correspondence therefore give the same transformation.

We could also, if $x, y, z$ are determined by (1), have put,
in place of (4),

\weq{6}{\xi=f\left(\dfrac{c+\partial\omega}{a}\right),\ \eta=\left(\dfrac{2}{\partial}\right)f_2\left(\dfrac{c+\partial\omega}{a}\right),\ \zeta=\left(\dfrac{2}{a}\right)f_1\left(\dfrac{c+\partial\omega}{a}\right)}

and consequently for $A, B$:

\weq{7}{\begin{aligned} 2A &= x\xi+(-1)^{\frac{n+1}{8}}(y\zeta+z\xi),\\ B &= x\xi y\zeta+x\xi z\eta+(-1)^{\frac{n+1}{8}}y\eta z\zeta \end{aligned}} % [sic] German source prints z\xi in 2A; the symmetric pattern suggests z\eta.

The fundamental curve and the degree of the correspondence
would then have remained the same, but the coincidences would have
acquired a different and indeed simpler meaning. In
both cases the coincidences belong to the singular values
of the modular functions, in which $\omega$ is an imaginary quadratic
irrationality, as we shall see more precisely in
what follows.

In the case $n\equiv 3 \pmod 8$ we set
\[
x=f^2(\omega),\ y=f_1^2(\omega),\ z=f_2^2(\omega),
\]
\weq{8}{\xi=f^2\left(\dfrac{c+\partial\omega}{a}\right),\ \eta=f_1^2\left(\dfrac{c+\partial\omega}{a}\right),\ \varrho=f_2^2\left(\dfrac{c+\partial\omega}{a}\right),}
%CONT>>
